% Supplementary source context; citations refer to the paper bibliography.
% This is a companion note, read with the paper, not a standalone proof.
Here $L(x)=x/\log x$ for $x>1$, $N_p$ is the first index whose
value is divisible by $p$, and $O$ counts odd non-tick used values
above the frontier in the state under discussion, as in the paper.
The notation and estimates are collected in the paper's Section~2.

The EKG sequence was discovered by Jonathan Ayres, who entered it in
the OEIS in September 2001 \cite{oeis}, and was studied by
Lagarias, Rains, and Sloane \cite{lrs}, who credit Ayres with its
discovery. It starts with $1,2$ and then chooses the least unused positive
integer having a common factor with the preceding term. It is
\href{https://oeis.org/A064413}{A064413} in the On-Line Encyclopedia
of Integer Sequences \cite{oeis}. Write $\U_t=\{a(1),\ldots,a(t)\}$.
The defining rule is equation~(1) of the paper.
Its first terms are $1,2,4,6,3,9,12,8,10,5,15,18,14,7,21$.
All LRS conjecture numbers in this edition refer to arXiv version~1.
Lagarias, Rains, and Sloane prove the permutation property, increasing
order of the primes (their Lemma 2), and linear upper and lower bounds.
Their numerical experiments suggested three branches: ordinary values
near $n$, primes near $n/2$, and three times primes near $3n/2$, as Ayres
had also observed \cite[Section~4]{lrs}.
Their PDF abstract says, exactly, ``We conjecture that almost all
$a(n)$ satisfy the asymptotic formula'' and displays
\[
 a(n)=n(1+1/(3\log n))+o(n/\log n)\quad\text{as }n\to\infty.
\]
Their Conjecture 2 defines $\sim$ as a ratio tending to $1$ and gives
these exact mathematical expressions, with the indicated conditions:
\[
\begin{array}{ll}
 a(n)\sim n[1+1/(3\log n)],&a(n)\ne p,3p\ (p\text{ prime}),\\
 a(n)\sim \tfrac12 n[1+1/(3\log n)],&a(n)=p\ (p\text{ prime}),\\
 a(n)\sim \tfrac32 n[1+1/(3\log n)],&a(n)=3p\ (p\text{ prime}).
\end{array}
\]
Thus the displayed Conjecture 2, read literally, asks only for three
ratio limits. We interpret its intended second-order content in the
stronger additive sense stated in their PDF abstract. The arXiv web
abstract has different parentheses; the PDF is the source here.

Hofman and Pilipczuk \cite[Theorem 4.3]{hp} prove that every odd prime
$p$ is preceded by $2p$; the following $3p$ is the consequence already
noted in LRS Conjecture 1. Their proof uses the LRS computation for
$p\leq25000$; Haver's proof reduces this finite range \cite{haver}.
Hofman and Pilipczuk also prove the three ratio limits in LRS Conjecture 2:
their Theorem 4.17 gives $a'_n=n+O(n/\log\log\log n)$.
Their $a'_n$ also replaces the prime $2$, so it agrees with our $b(n)$
only up to finitely many indices and yields the same asymptotic bound.
Their Conjecture 5.1 writes the LRS formula with $\sim$, but they
introduce it as ``the stronger conjecture''; we read this as support for
the additive second-order reading above.
Haver \cite{haver} proves the
stronger $b(n)=n+O(n/\log n)$ and quantitative estimates for early even
values, used values above the least unused even integer, and missing
composites below it. Haver's Corollary 2(e), in version 3 dated October 4,
2026, reduces the second-order assertion at these even integers to a
negligible count of odd non-tick used values above them. Haver's Theorem 1(d)
(also stated in Haver's abstract) and Corollary 12(d) already give
$E(t)-t\leq(1/3+o(1))t/\log t$ and
$N_p\geq2p-(2/3+o(1))p/\log p$, with the conjectured coefficients, whose sharpness follows from the matching asymptotic.
Our contribution is the lower half of the clock,
$E(t)-t\geq(1/3-o(1))t/\log t$, and the transfer from frontier times
to the full sequence of indices.
We prove the remaining hypothesis in Haver's Corollary 2(e), and a
uniform bound for missing composites needed for that transfer.
The reduction and the estimates it uses are Haver's.

An explicit rate looks within reach using Brun--Titchmarsh in place
of the fixed-modulus sieve; this is left for future work.

\section{Related work}\label{sec:related}

The argument completes the second-order program of Haver \cite{haver},
using the structural work of Lagarias, Rains, and Sloane \cite{lrs} and
Hofman and Pilipczuk \cite{hp}. Haver's Theorem 1(d) and Corollary 12(d)
already prove the upper half of the frontier clock with sharp constant
$1/3$ in Haver's manuscript; Haver's Lean sources do not cover the sharp
coefficient or the Corollary 2(e) criterion. The new contribution is the
lower half plus the transfer to the full sequence of indices.
Haver also proves the exact counting identity and the equivalence stated
in Corollary 2(e), and identifies the rough-predecessor obstruction in
the remark after Corollary 12. We discharge the remaining hypothesis of
that equivalence and prove the uniform missing-composite estimate needed
for the transfer, extending the Lean proof to the full second-order theorem.
The prime-capacity lemma is an elementary fixed-modulus sieve bound
\cite{sievebook}, proved here; it does not require a short-interval
prime number theorem. Haver's cited manuscript is unrefereed.


The reference to the sieve text supplies background only. The interval
prime-capacity lemma is proved in full here; no estimate from that book is
an imported hypothesis or an unchecked proof step.

\section{Earlier source comparison}
Lagarias, Rains and Sloane's arXiv v1 PDF \cite[abstract]{lrs} says
``We conjecture that almost all $a(n)$ satisfy the asymptotic formula''
and displays $a(n)=n(1+1/(3\log n))+o(n/\log n)$.
Their opening credits the sequence as ``discovered by Jonathan Ayres'';
Section~4 says of the three lines, ``This was also observed by Ayres''.
Their arXiv Conjecture~2 uses ratio
asymptotics, a weaker assertion. Hofman and Pilipczuk's smoothed sequence $a'_n$ replaces $p,3p$
by $2p$ for every prime $p$, including $2$. Their Theorem~4.17 \cite{hp} gives $a'_n=n+O(n/\log\log\log n)$.
This convention differs from $b$ at only finitely many indices.
Haver's version~3 \cite[Corollary~2(e)]{haver} states the counting
criterion; Haver's Theorem~1(d) and Corollary~12(d) give the sharp upper
bound $E(t)-t\leq(1/3+o(1))t/\log t$, where $E(t)$ is the least
unused positive even integer after $t$ terms. Here the new work is the
remaining lower bound and the transfer to all smoothed indices.
The reduction and the estimates it uses are Haver's.
The hard step is improving the displacement estimates from
$O(x/\log x)$ to $o(x/\log x)$; multiplying the existing bounds
leaves a fixed coefficient. The proof below supplies sparsity of the
prime cofactors that those bounds alone do not detect.


\subsection*{Further comparisons and conventions}
These are the three lines observed by Ayres and by Lagarias, Rains, and
Sloane \cite[Section~4]{lrs}. For large $k$ with $a(k)\in\{p,3p\}$ we have $b(k)=2p\leq2k$, so
at most $2\pi(N)+O(1)=o(N)$ indices $k\leq N$ are of these two kinds.
Hence the first line gives the formula of the LRS abstract for almost
all $n$: this is the additive second-order interpretation. It implies
the literal Conjecture 2, whose ratio limits were already proved
by Hofman and Pilipczuk \cite[Theorem~4.17]{hp}. LRS Conjecture 3,
$13n/28\leq a(n)\leq12n/7$ for every $n$, with equality only at $n=28$
and $n=7$, respectively, is not settled here: an eventual asymptotic
does not settle the required finite range. It is proved, with computer
assistance, in our companion paper \cite[Theorem~1, version~1]{ekgbounds}. Nor do we address the
unproved distributional suggestions for higher-order fluctuations in
LRS Section 4 or their unnumbered cycle questions, beyond the bounded
cycle lengths of Section~\ref{sec:consequences}.

Section~\ref{sec:consequences} derives three consequences. For all
large $n$, $a(n)>n$ exactly when $a(n)$ is composite, and
$\#\{k\leq N:a(k)<k\}\sim N/(2\log N)$, so
$\#\{k\leq N:a(k)>k\}/N\to1$. This proves two 2009 conjectures of
Thomas Ordowski recorded in OEIS A064413: that $a(n)>n$ almost
everywhere when $a(n)$ is not prime, which our eventual form
strengthens, and the density statement. The equations
$a(k)=k$, $2a(k)=k$ and $2a(k)=3k$ have finitely many solutions; this
proves Artur Jasinski's 2008 conjectures that A152458 ($a(k)=k$) is
finite and, on A152519, that the three lines carry finitely many points
in total. And the permutation has finitely many cycles of each bounded
length.

Hofman and Pilipczuk's smoothed sequence $a'_n$ replaces $p,3p$
by $2p$ for every prime $p$, including $2$. Their Theorem~4.17 \cite{hp} gives $a'_n=n+O(n/\log\log\log n)$.
This convention differs from $b$ at only finitely many indices.


We use these established estimates for the greedy sequence. The prime pattern goes back to
Hofman and Pilipczuk \cite[Theorem~4.3]{hp}; we use its introduction formulation
as proved in Haver's archive \cite[Section~2, (F2)]{haver}.
The early-even bound follows from Haver's bound
$K(t)\leq12\pi(\sqrt e)+720$ \cite[version~3, Theorem~1(a)]{haver}
and elementary prime estimates.
One may take $C_1=4\cdot793^2+2\cdot793$ in \eqref{eq:oldspan}.
The printed sources of \eqref{eq:oldhole} and \eqref{eq:oldrate}
are Haver's Theorem 11 and Corollary 12(c), respectively.
The constants in \eqref{eq:early}--\eqref{eq:F} are convenient
weakenings proved in the accompanying sources, rather than optimized
constants: $793=12\cdot6+721$ and $802=13\cdot6+724$ follow from the
eventual bound $\pi(\sqrt e)\leq6\sqrt e/\log e$ and
$\sqrt e/\log e\geq1$, absorbing the additive constants.
The prime estimate in \eqref{eq:F} does not require the
prime number theorem.

For the following comparison, write $e=E(t)$, let $D$ denote a positive
displacement, let $r$ be the least prime factor of a predecessor, and let
$h$ denote a missing composite. The quantities $k(t),K_0,F(e)$ and
$\spf$ have been defined above.

The coefficient has the counting heuristic of LRS Section 4 \cite{lrs}:
\[
 n\approx x-\frac{x}{2\log x}+\frac{x}{6\log x}
       =x-\frac{x}{3\log x}.
\]
Our new step upgrades two $O(x/\log x)$ displacement bounds to
little-$o$. For the upper bound, the old greedy inequality is
$D\leq2rk(t)$, with $r\leq\sqrt{e+D}$ and
$k(t)\leq K_0\sqrt e/\log e$. Multiplying these estimates leaves a
fixed coefficient of $e/\log e$; it does not force that coefficient
to vanish. For the lower bound,
$D\leq2\spf(h)F(e)$ and $\spf(h)\leq\sqrt e$ have the same loss.
Haver already identifies these rough-predecessor and rough-missing-composite
obstructions in the remark after Corollary 12 \cite{haver}.
The additional information is sparsity among the forced prime cofactors. Given a putative displacement $D\geq\delta x/\log x$, we
instead select multiples $2rMj$ for a fixed auxiliary integer $M$ chosen
after $\delta$ and positive integers $j$ and force their distinct neighbors
to be products $rq_j$ with $q_j$ prime. There are at least
$D/(4rM)-O(1)$ such cofactors in an interval of length
$O(x/(r\log x))$. For each fixed $\delta>0$ their density is bounded
below by a positive constant independent of $x$. It need not be
independent of $\delta$: $M$ also depends on $\delta$. The sieve bound
permits an arbitrarily small fixed upper density, chosen after $\delta$
and $M$, which is exactly what the little-$o$ conclusion needs.

The two estimates are an upper bound for the largest ordinary used
value above the least unused even integer, and a lower bound on every
unused composite below it. Sections~\ref{sec:upper} and
\ref{sec:lower} prove them, then exact counting and the prime number
theorem give the result. We extract no explicit remainder rate. No explicit cutoff $N(\eps)$ results.
The nonquantitative step is the choice of a fixed sieve modulus for
each $\delta$, followed by unspecified eventual cutoffs in localization
and the qualitative prime number theorem. This is not a claim of an
intrinsically ineffective obstruction: the sieve constants are
computable, but this proof does not supply a function $N(\eps)$.

\subsection*{Original numerical display}


Table~\ref{tab:ratio-original} gives the normalized error
$R(n)=(b(n)-n)\log n/n$, whose limit is $1/3$ by
Theorem~\ref{thm:main}. These values illustrate the formula; they are
not used in its proof and do not imply a monotone approach to the limit.
\begin{table}[ht]
\centering
\begin{tabular}{r r c c}
$n$ & $b(n)$ & $R(n)$ & $1/3$\\\hline
1,000 & 1,043 & 0.297033477 & 0.333333333 \\
10,000 & 10,361 & 0.332493287 & 0.333333333 \\
100,000 & 103,088 & 0.355519138 & 0.333333333 \\
1,000,000 & 1,024,686 & 0.341049694 & 0.333333333 \\
10,000,000 & 10,210,804 & 0.339775904 & 0.333333333 \\
100,000,000 & 101,839,957 & 0.338932605 & 0.333333333 \\
500,000,000 & 508,457,544 & 0.338811220 & 0.333333333 \\
\end{tabular}
\caption{The smoothed EKG error, with absolute rounding error at most
$5\cdot10^{-10}$ in each displayed ratio.}\label{tab:ratio-original}
\end{table}


\subsection*{Schematic critical lines}
\begin{figure}[ht]
\centering
\begin{tikzpicture}[x=1.1cm,y=0.65cm]
\draw[->] (0,0) -- (5.1,0) node[right] {$n$};
\draw[->] (0,0) -- (0,7.7) node[above] {$a(n)$};
\draw[blue] (0,0) -- (4.5,2.25) node[right,black] {$n/2$};
\draw[blue] (0,0) -- (4.5,4.5) node[right,black] {$n$};
\draw[blue] (0,0) -- (4.5,6.75) node[right,black] {$3n/2$};
\end{tikzpicture}
\caption{The three critical lines, shown schematically. This is not
sequence data; Corollary~\ref{cor:branches} gives the positive
second-order corrections, and Corollary~\ref{cor:lines} proves
finiteness of intersections with these exact lines.}
\label{fig:branches-schematic}
\end{figure}
\subsection*{Consequences and supplementary explanations}
Ordowski's first conjecture asks only that $a(n)>n$ ``almost
everywhere'' when $a(n)$ is not prime. The first statement, for all but
finitely many $n$, is our stronger form; it implies Ordowski's conjecture in
the density-one reading of ``almost everywhere''. The last statement is Ordowski's second
conjecture. The count of prime values alone,
$\#\{k\leq N:a(k)\text{ prime}\}\sim N/(2\log N)$, already follows from
the ratio form; the new content is that these are, with finitely many
exceptions, exactly the indices below the diagonal.

Two related OEIS entries are also finite sets: Ordowski's ``exceptional
numbers'' A155549, the values $a(n)$ that are not prime and satisfy
$a(n)\leq n-1$ (its terms $16,26,49,\ldots$ occur at the indices
$n=17,27,50,\ldots$), and A064420, the $k$ with $a(k+1)=k$. Both are
listed as ``probably complete''. By Corollary~\ref{cor:sign} the indices
$n$ with $a(n)$ not prime and $a(n)\leq n-1$ form a finite set, so its
image A155549 is finite. For the
second, $a(n)=n-1$ fails for large $n$: a prime value has
$a(n)\leq2n/3$ and any other value has $a(n)>n$. De Vlieger
\cite{devlieger} remarked that it is ``not surprising'' that A064420
is finite.



Lagarias, Rains and Sloane \cite[Section~7]{lrs} note that the
permutation appears to have infinitely many finite cycles and
infinitely many infinite cycles, and expect both to be hard to prove.
The next result implies that infinitely many finite cycles, if they
exist, have unbounded lengths.

The letters $M$ and $L$ used locally for an integer
or a cycle length are independent of the function $L(x)$ and of each
other; $\N$ always denotes the natural numbers.

Haver identifies these obstructions in the remark after Corollary~12
\cite{haver}. The cofactor argument and the order of its constants
are described in the comparison above.

We first bound the used ordinary values above $E(t)$, then the unused
composites below it. The finite-prefix argument used below is the
maximum-of-initial-values argument preceding \eqref{eq:adjacent}.
The terms frontier before and after a value mean $E(j-1)$ and $E(j)$
at its occurrence index $j$; the upper-excursion cutoff is the index
threshold supplied by the first proposition.

The clock estimates in the following proof refer to \eqref{eq:clock};
finite seed exclusions are the finitely many initial indices omitted by
the eventual estimates. The two excursion propositions are
Propositions~\ref{prop:upper} and \ref{prop:lower}.

The formalized sequence is the canonical
greedy sequence of \eqref{eq:rule}, shifted from zero-based to one-based
indexing, with \eqref{eq:smooth} applied only to odd primes and their
triples; $2$ and $6$ stay fixed. At $n=0$ it sets $a(0)=1$, outside the
stated range. The inequality is over the reals with natural logarithms.

Lean \cite{lean}, mathlib \cite{mathlib} and the prime number theorem
library \cite{pntcode} support the formal proof; the literal definitions
and source correspondence are in the accompanying checking notes.

Lagarias, Rains and Sloane credit Ayres with the discovery and the
three lines \cite[Section~1, first paragraph; Section~4]{lrs}; their abstract conjectures
the displayed second-order formula at almost all indices
\cite[abstract]{lrs}. Their Conjecture~2 instead asks for ratio limits,
proved by Hofman and Pilipczuk \cite[Theorem~4.17]{hp}.

\subsection*{Additional source and estimate details}
Consequences
include an eventual strengthening of Ordowski's inequality conjecture
and Ordowski's density-one conjecture, Jasinski's finiteness conjectures on the
three critical lines, the below-diagonal counting asymptotic and finiteness
of cycles of each bounded length.

The proof is formalized in Lean, as described in Section~\ref{sec:lean};
the three-branch formulas and the sharper complementary count are written
deductions from the checked results.

The constants $793=12\cdot6+721$ and $802=13\cdot6+724$ are convenient
weakenings using $\pi(\sqrt e)\leq6\sqrt e/\log e$ and
$\sqrt e/\log e\geq1$ eventually; $C_1=4\cdot793^2+2\cdot793$ suffices.

Equation~\eqref{eq:oldrate} also bounds adjacent ordinary values.
If $y=a(j)$ is unchanged by smoothing and $e/2<y<2e$, then, uniformly
as $e\to\infty$, the occurrence index $j$ tends to infinity.
To see this, choose any fixed index cutoff and exceed the maximum
of the finitely many preceding values. Equation~\eqref{eq:oldrate}
then gives $j/2\leq y\leq3j/2$ and hence $e/3<j<4e$.
If an adjacent value $d=a(j-1)$ or $a(j+1)$ is also unchanged by
smoothing, \eqref{eq:oldrate} at both indices gives
\begin{equation}\label{eq:adjacent-notes}
 |d-y|\leq C L(j)+C L(j\pm1)+1=O(L(e)).
\end{equation}
The constant is uniform. Here the constant depends only on \eqref{eq:oldrate}.



The bound $O=o(L(e))$ proves the remaining hypothesis of the equivalence
stated in Haver's Corollary 2(e) and proved in Corollary 12(d);
the counting criterion itself is not new here.

The proof uses Lean \cite{lean}, mathlib \cite{mathlib} and the prime
number theorem library \cite{pntcode}; the checking notes give the source correspondence.

\subsection*{Convention and consequence comments}
The smoothing of Hofman and Pilipczuk also replaces $2$, and differs
from ours at only finitely many indices \cite[Theorem~4.17]{hp}.

Corollary~\ref{cor:sign} also implies finiteness of Ordowski's exceptional
values A155549 and the indices A064420. The checking notes give the
short deductions and De Vlieger's ``not surprising'' observation
about A064420 \cite{devlieger}.

The bounded-cycle question comes from LRS \cite[Section~7]{lrs}.
The next result forces any infinite collection of finite cycles to
have unbounded lengths.

The Lean theorem \texttt{ekg\_sharp\_all\_index} states that, for every
real $\eps>0$, there is a natural $N\geq3$ such that every natural
$n\geq N$ satisfies \eqref{eq:epsilon}.

Haver's archive supplies the general counting identities
\cite[Corollary~12(b)]{haver}; our formalization establishes their
prime-interval and tick endpoint hypotheses and derives the coefficient
and transfer to all indices.

In the proofs, the old rate and old span mean \eqref{eq:oldrate} and
\eqref{eq:oldspan}, respectively. The new step replaces the two
$O(L(e))$ displacement bounds by $o(L(e))$ bounds.

The interval-prime
bound below gives the contradiction. The fixed sieve modulus and eventual
cutoffs yield no explicit function $N(\eps)$ in this proof.

All logarithms are natural; $\N$ denotes the natural numbers.

\subsection*{Additional conjecture context}
Theorem~\ref{thm:main} has the following consequences for questions
recorded in the OEIS \cite{oeis}. On A064413, Thomas Ordowski
conjectured on January 23, 2009, ``If $a(n)\neq p$, then almost
everywhere $a(n)>n$'' and ``$\lim\#(a_n>n)/n=1$''. Artur Jasinski, in
December 2008, called $a(k)=k$, $2a(k)=3k$ and $2a(k)=k$ the three
critical lines of the sequence. Jasinski's entry A152458, the indices with
$a(k)=k$, records the conjecture that it is finite, and on A152519 (the
indices with $2a(k)=3k$) Jasinski conjectured that the number of points on the
three lines is finite, asking also whether the listed points are
complete. Below, $p$ denotes a prime.

Amiram Eldar recorded in 2024 that neither sequence has a
further term below $10^5$ \cite{oeis}. These computations are distinct
from the eventual finiteness result.

All LRS section, figure and conjecture numbers used here and in the
article follow arXiv v1, identified in the unchanged bibliography.
Haver's Theorem~1(b), with Theorem~8(b), gives the non-tick span estimate
\eqref{eq:oldspan}. Haver's Corollary~12(b) gives the counting identity in
prime-interval form; the article displays the corresponding prime-count
form and its boundary corrections.
